\begin{resumo} 
A geração de energia fotovoltaica varia conforme a disponibilidade de radiação solar, as condições meteorológicas e as características da instalação. Este trabalho avalia o uso de redes neurais recorrentes do tipo Long Short-Term Memory (LSTM) para estimar a energia gerada na próxima hora. Foi desenvolvida uma prova de conceito para um sistema hipotético de 5,5 kWp, com inversor de 5 kW, localizado em Recife, Pernambuco. Os experimentos utilizaram 120 dias de dados sintéticos, com perturbações meteorológicas simuladas, e compararam o LSTM direto, o LSTM residual, a RNN simples, o Random Forest e métodos de referência. O LSTM direto com perda de Huber apresentou o menor erro no teste, com MAE de 0,0572 kWh, WAPE de 6,89\% e coeficiente de determinação de 0,9910. O LSTM residual reduziu em 28,47\% o MAE da referência física. Os resultados demonstram o potencial da abordagem como protótipo acadêmico.
\\\\
\noindent\textbf{Palavras-chave:} energia solar fotovoltaica; aprendizado de máquina; séries temporais; LSTM; previsão meteorológica.

\end{resumo}